\chapter{Risultati e conclusioni}

Al termine della fase di sviluppo è stato eseguito un test automatizzato con l’obiettivo di misurare la qualità complessiva del sistema e raccogliere metriche quantitative riconducibili ai principali KPI individuati e descritti di seguito. Questa sezione esamina inoltre i limiti dell'implementazione corrente e le possibili linee di sviluppo futuro, ponendo le basi per le riflessioni conclusive.

\section{Testing automatizzato}

La valutazione del sistema è stata predisposta tramite uno script automatizzato contenuto nella cartella \textit{eval}. In particolare, il file \textit{run\_kpi\_eval.py} esegue una serie di test simulati a partire dal dataset \textit{test\_dataset.jsonl} e produce un riepilogo dei risultati nel file \textit{eval/outputs/results.csv}.

Il dataset utilizzato per il testing non rappresenta interazioni reali raccolte da utenti esterni, ma un insieme di casi d’esempio costruiti per verificare il comportamento del sistema rispetto a domande il più possibile vicine a quelle che potrebbero essere formulate dagli utenti finali. I casi sono stati definiti cercando di mantenere eterogeneità nei contenuti, nei domini informativi e nelle possibili modalità di richiesta. Di conseguenza, i risultati ottenuti misurano la qualità del sistema in un ambiente controllato, ma non certificano in modo definitivo le prestazioni in uno scenario reale, caratterizzato da utenti esterni, formulazioni imprevedibili e richieste ambigue.

Il dataset di test contiene \textit{130} casi testuali. Le domande coprono diversi domini informativi, tra cui informazioni generali, venue, calendario, ticketing, volontariato e accessibilità. La maggior parte dei casi è in italiano, ma sono presenti anche esempi in inglese, spagnolo, francese e arabo, così da verificare il comportamento multilingua del sistema. Per ogni caso vengono specificati l’identificativo del test, la modalità, il messaggio, la lingua, il dominio atteso, i documenti rilevanti attesi, il feedback simulato e, in alcuni casi, l’eventuale apertura di un ticket.

Il runner esegue prima una fase di validazione del dataset. In questa fase controlla che ogni caso abbia un identificativo, che non ci siano duplicati, che la modalità sia ammessa e che i campi richiesti siano coerenti con il tipo di test. Ad esempio, un caso testuale deve avere un messaggio, mentre eventuali casi audio o immagine dovrebbero indicare un file associato.

I KPI individuati possono essere organizzati in tre macroaree. La prima è la qualità informativa, misurata da \textit{Domain Accuracy}, \textit{Recall@5}, \textit{Precision@5}, \textit{Mean Reciprocal Rank} (MRR) e \textit{Source Coverage Rate}. La seconda è la performance tecnica, misurata da \textit{Average Latency}, \textit{p95 Latency} ed \textit{Error Rate}. La terza è l’impatto operativo, misurato da \textit{Escalation Rate}, \textit{Containment Rate} e \textit{Feedback Score}. La qualità informativa dipende dalla knowledge base e dal retrieval, la qualità conversazionale dipende dal frontend, dal backend e dal modello linguistico, mentre la dimensione operativa dipende da ticket, dashboard e feedback.

Per misurare la qualità informativa, il runner importa direttamente le funzioni di query planning e retrieval del backend, costruisce il piano della richiesta, interroga il servizio RAG e confronta i documenti recuperati con quelli attesi nel dataset. Per quanto riguarda la performance tecnica, invece, il runner invia una richiesta all’endpoint \textit{/api/chat}, costruendo un \textit{session\_id} specifico. Il backend elabora la richiesta come avverrebbe in una normale conversazione, genera la risposta, restituisce eventuali fonti, salva i messaggi e produce le informazioni necessarie per feedback e ticket. Il runner misura quindi latenza, successo della chiamata, presenza di errori, copertura delle fonti ed eventuale comportamento di escalation.

Il flusso prevede anche la possibilità di pubblicare feedback e ticket simulati. Se l’opzione è abilitata, il runner invia il feedback all’endpoint \textit{/api/messages/message\_id/feedback}. Nei casi predisposti per l’escalation, invia inoltre una richiesta all’endpoint \textit{/api/tickets}.

Il runner confronta lo stato del database e calcola metriche come escalation rate cumulativo, containment rate cumulativo, feedback score e delta dei nuovi ticket rispetto alle nuove conversazioni create durante la run. Questa scelta collega la valutazione automatica ai dati realmente salvati dal sistema.

Il file \textit{results.csv} generato dal runner contiene un riepilogo aggregato dei risultati ottenuti. Ogni riga del \textit{CSV} riporta il nome del KPI, una breve descrizione, il valore calcolato e l’unità di misura.

\section{Valutazione metriche e KPI}

Pur in assenza di un \textit{Service Level Agreement} (SLA) formalmente definito, i risultati emersi dal file CSV evidenziano comunque un comportamento complessivamente positivo del sistema, almeno con riferimento a uno scenario controllato basato su richieste utente di base. Le metriche di qualità informativa indicano che il sistema riesce, nella maggior parte dei casi, a classificare correttamente il dominio della richiesta e a recuperare documenti pertinenti dalla knowledge base. Le metriche operative evidenziano invece una buona capacità di contenere le richieste informative senza ricorrere sistematicamente all’intervento dell’operatore umano.

La Tabella~\ref{tab:kpi-results} riporta i principali KPI registrati nel file \textit{results.csv}.

\begingroup
\small
\begin{longtable}{p{0.24\textwidth}p{0.18\textwidth}p{0.48\textwidth}}
\caption{Risultati ottenuti dal test automatizzato.}
\label{tab:kpi-results}\tabularnewline
\toprule
\textbf{KPI} & \textbf{Descrizione} & \textbf{Interpretazione}\tabularnewline
\midrule
\textbf{Domain Accuracy} & 92.62\% & Il sistema classifica correttamente la maggior parte dei domini informativi delle domande di test.\tabularnewline
\textbf{Recall@5} & 79.51\% & In circa quattro casi su cinque almeno un documento rilevante compare tra i primi cinque risultati recuperati.\tabularnewline
\textbf{Precision@5} & 74.75\% & La quota media di documenti pertinenti nei primi cinque risultati è elevata, indicando un retrieval complessivamente selettivo.\tabularnewline
\textbf{MRR} & 0.7204 & Il primo documento corretto tende a comparire in posizione alta quando viene recuperato.\tabularnewline
\textbf{Source Coverage Rate} & 79.07\% & La maggior parte delle risposte valide contiene almeno una fonte recuperata dalla knowledge base.\tabularnewline
\textbf{Average Latency} & 5.65 s & Il tempo medio end-to-end è basso.\tabularnewline
\textbf{p95 Latency} & 10.92 s & Il 95\% delle richieste rientra entro circa 11 secondi.\tabularnewline
\textbf{Error Rate} & 0.77\% & Le richieste fallite sono molto poche rispetto al totale dei casi valutati.\tabularnewline
\textbf{Escalation Rate} & 6.20\% & Una quota limitata di conversazioni simulate viene trasferita o trasformata in ticket operatore.\tabularnewline
\textbf{Containment Rate} & 93.80\% & La maggior parte delle richieste simulate viene gestita senza intervento umano.\tabularnewline
\textbf{Feedback Score} & 92.73\% & I messaggi valutati nel database risultano prevalentemente soddisfacenti.\tabularnewline
\bottomrule
\end{longtable}
\endgroup

La Domain Accuracy del 92.62\% è uno dei risultati più solidi. Indica che il sistema riesce a riconoscere correttamente il dominio della richiesta nella maggior parte dei casi. Questo risultato è importante perché il dominio influenza il retrieval e la costruzione della risposta. Se il dominio è errato, il sistema rischia di interrogare documenti meno pertinenti e di fornire un contesto meno adatto al modello linguistico.

Recall@5 pari al 79.51\% indica che, nella maggior parte dei casi, almeno un documento rilevante entra tra i primi cinque risultati recuperati. In un sistema RAG questo aspetto è fondamentale, perché il modello linguistico può generare una risposta fondata solo se il documento corretto entra nel contesto. Quando il documento rilevante non viene recuperato, anche un buon modello linguistico può rispondere in modo incompleto o troppo generico.

Precision@5 pari al 74.75\% è un risultato positivo, perché indica che una parte consistente dei primi cinque documenti recuperati è pertinente rispetto alla domanda.

Con un MRR del 72\%, i dati confermano che i documenti pertinenti recuperati compaiono tipicamente nei primi risultati. Questo è rilevante perché i primi documenti influenzano maggiormente il prompt e quindi la risposta finale. Un MRR vicino a 1 indicherebbe che il primo risultato è quasi sempre corretto.

Il Source Coverage Rate del 79.07\% indica che la maggior parte delle risposte valide viene accompagnata da almeno una fonte. Questo dato è importante in un sistema informativo, perché le fonti aumentano tracciabilità e fiducia.

Il tempo medio di risposta è di circa 6 secondi, mentre il p95 arriva a circa 11 secondi. Questi valori sono accettabili per un servizio conversazionale in tempo reale, dove l’utente si aspetta risposte quasi immediate.

L’Error Rate pari allo 0.77\% è molto basso e indica che il sistema fallisce raramente durante la valutazione. Questo dato è particolarmente significativo, poiché riflette la buona affidabilità del flusso applicativo nel suo complesso.

Dal punto di vista operativo, Escalation Rate e Containment Rate devono essere letti insieme. L’Escalation Rate del 6.20\% indica che una parte limitata delle conversazioni simulate viene trasferita o trasformata in ticket. Il Containment Rate del 93.80\% indica invece che la maggior parte delle richieste di test viene gestita senza intervento umano.

Il Feedback Score del 92.73\% è positivo e rispecchia la soddisfazione degli utenti in riferimento all'intero flusso di utilizzo del sistema.

\section{Limiti e sviluppi futuri}

Lo stato attuale del progetto, pur mostrando un pieno funzionamento del prototipo, presenta alcuni limiti e criticità che devono essere considerati in vista di una possibile evoluzione verso un utilizzo reale.

Un primo limite riguarda la natura simulata del benchmark. Il dataset utilizzato rende la valutazione ripetibile e controllabile, ma non può rappresentare pienamente il comportamento degli utenti che potrebbero utilizzare l’applicazione durante l’evento. In un contesto operativo reale, è opportuno considerare che le richieste degli utenti possono assumere forme notevolmente più complesse e articolate rispetto a uno scenario sperimentale. Esse possono infatti contenere errori ortografici o sintattici, ambiguità semantiche, formulazioni incomplete o ellittiche, nonché riferimenti impliciti a contesti non dichiarati. Tali fattori introducono un livello di variabilità e incertezza che mette alla prova le capacità di comprensione e adattamento dell'agente AI, richiedendo soluzioni robuste in grado di gestire la naturale imprecisione del linguaggio umano senza compromettere l'affidabilità della risposta.

Un secondo limite riguarda la valutazione multimodale. Il sistema supporta input audio e immagini, ma queste modalità non sono ancora incluse nel benchmark automatico. Per una valutazione più completa sarebbe necessario costruire un dataset dedicato, composto da file audio, immagini reali, screenshot, foto sfocate e contenuti acquisiti in condizioni non ideali. Sarebbero inoltre necessarie metriche specifiche per valutare la qualità della trascrizione audio, dell’OCR e della descrizione visuale.

Un ulteriore limite riguarda la gestione delle informazioni dinamiche. La knowledge base è adatta a contenere informazioni relativamente stabili, ma non può sostituire fonti live o dati soggetti ad aggiornamenti frequenti, che richiederebbero integrazioni con API ufficiali o sorgenti aggiornate automaticamente. Quando i dati non sono disponibili, il sistema deve adottare un comportamento chiaro e prevedibile: segnalare all'utente l'indisponibilità dell'informazione, senza alcun tentativo di supplire con risposte generate in assenza di fonti. In termini operativi, questo significa rispondere con un messaggio standardizzato che spieghi i motivi della mancata risposta, evitando qualsiasi forma di "invenzione" o di completamento speculativo e orientando eventualmente l'utente verso canali alternativi (operatore umano, documentazione, FAQ).
Questa policy non è un limite, ma una garanzia, in quanto tutela la qualità del servizio e la credibilità del sistema, evitando che un errore informativo si trasformi in un danno per l'utente o per l'azienda.

Anche l'aspetto della sicurezza offre margini di intervento e ottimizzazione. Il sistema include già alcuni elementi di sicurezza di base, ma una versione di produzione richiederebbe interventi più strutturati, come la gestione dei ruoli, l’uso stabile di HTTPS, una configurazione più restrittiva del \textit{CORS}, meccanismi di \textit{rate limiting}, monitoraggio, backup, protezione degli upload e gestione conforme dei dati personali nel rispetto del GDPR e delle procedure tipiche di un servizio di customer care.

Un possibile sviluppo riguarda lo smistamento dei ticket tra più operatori. Nello stato attuale, la dashboard dimostra il passaggio dalla chat al ticket, ma non implementa ancora un sistema completo di assegnazione automatica. In un contesto operativo reale, la distribuzione dei ticket deve essere governata da una pluralità di criteri interconnessi, che riflettono la complessità e la variabilità delle richieste in ingresso. Tra questi figurano il dominio di appartenenza della richiesta, il livello di priorità assegnato, la lingua dell'utente, la reperibilità e le competenze specifiche degli operatori, nonché l'eventuale carico di lavoro in essere. Una corretta parametrizzazione e integrazione di tali fattori è condizione necessaria per un'allocazione efficiente delle risorse e per il mantenimento di standard qualitativi elevati nel servizio di customer care.

Anche il perimetro di intervento degli operatori potrebbe essere ampliato. Attualmente l’operatore è gestito in modo statico tramite database, mentre una versione futura dovrebbe consentire all'amministratore del sistema di aggiungere nuovi operatori autorizzati, disabilitare gli account, modificare i ruoli e assegnare permessi differenziati.

Tra i possibili sviluppi futuri del sistema si annovera l'implementazione di un meccanismo di verifica dell'indirizzo e-mail dell'utente prima della conclusione della procedura di apertura del ticket. Allo stato attuale, la piattaforma si limita a una validazione formale dell'indirizzo senza tuttavia accertare che l'utente sia effettivamente in possesso delle credenziali di accesso alla casella di posta indicata. Tale limite può tradursi in criticità operative non trascurabili, quali l'impossibilità di recapitare comunicazioni successive relative al ticket, con conseguente insoddisfazione dell'utente e potenziale aggravio per il customer care. L'adozione di un sistema di conferma via e-mail, tramite link dedicato o codice monouso, rappresenterebbe una soluzione efficace per garantire la tracciabilità e la certezza del contatto, migliorando la qualità complessiva del servizio. L’introduzione di un flusso con codice \textit{One-Time Password} (OTP) inviato via e-mail permetterebbe di ridurre ticket falsi, errori di digitazione e richieste non contattabili.

Un'ultima area di intervento concerne il modulo di retrieval, ovvero il processo di recupero dei documenti pertinenti al fine di supportare la generazione delle risposte. I risultati finora conseguiti si rivelano complessivamente positivi, attestando la bontà dell'architettura implementata. Tuttavia, ulteriori affinamenti, in termini di indicizzazione semantica, pesatura dei termini e strategie di ranking, potrebbero contribuire a incrementare la precisione e la completezza del recupero, limitando il rischio di fornire contesti parziali, duplicati o solo parzialmente coerenti con la richiesta dell'utente. Il perfezionamento di questa componente si configura come un passo strategico per elevare la qualità complessiva del sistema e consolidarne l'affidabilità.

\section{Considerazioni finali}

In conclusione, il project work ha permesso di affrontare il ciclo di vita completo di una soluzione software innovativa, sviluppata a partire da un’esigenza concreta e collegata a un obiettivo di business ben definito. Il risultato ottenuto è un prototipo funzionante, documentato e valutato tramite metriche quantitative e qualitative, in grado di mostrare un comportamento complessivamente positivo rispetto agli scenari di test considerati.

Il sistema risponde al problema principale da cui è partita la progettazione, ovvero la frammentazione delle informazioni relative ai Giochi del Mediterraneo Taranto 2026. Attraverso un’unica interfaccia conversazionale, l’utente può infatti accedere a informazioni rilevanti sui diversi domini dell’evento, senza dover consultare manualmente fonti differenti o individuare in autonomia il canale corretto a cui rivolgersi.

La conclusione principale che si trae dal presente lavoro è che un agente di intelligenza artificiale destinato al customer care non può limitarsi alla generazione automatica di risposte testuali. Affinché possa costituire un effettivo valore aggiunto, è necessario che sia supportato da dati affidabili e costantemente aggiornati, da fonti chiare e verificabili, da una dichiarazione trasparente dei propri limiti informativi, da meccanismi di escalation ben definiti verso l'operatore umano e da indicatori di performance (KPI) in grado di valutarne l'efficacia complessiva, non solo in termini di rapidità, ma anche di accuratezza, appropriatezza e impatto sull'esperienza del cliente.

Un possibile scenario di sviluppo futuro può essere articolato in tre fasi. La prima riguarda il consolidamento del prototipo, con la correzione di eventuali bug residui, il miglioramento delle metriche, l’introduzione dello streaming della risposta, l’uso di meccanismi di \textit{caching} e il completamento dei test automatici. La seconda fase riguarda la validazione controllata con utenti reali o con dataset multimodali più ampi, includendo una gestione più avanzata degli operatori, la verifica dell’e-mail e ulteriori ottimizzazioni della pipeline RAG. La terza fase riguarda invece l’eventuale messa in produzione su infrastruttura cloud, con maggiore attenzione a sicurezza, privacy, monitoraggio, SLA e integrazione con API real-time ufficiali.

\begin{table}[H]
\centering
\caption{Scenario evolutivo del progetto.}
\label{tab:roadmap-evoluzione}
{\small
\begin{tabular}{p{0.22\textwidth}p{0.32\textwidth}p{0.36\textwidth}}
\toprule
\textbf{Fase} & \textbf{Obiettivo} & \textbf{Interventi}\tabularnewline
\midrule
\textbf{Consolidamento} & Rendere stabile il prototipo. & Bug fixing, streaming, caching, ottimizzazione dei prompt, test automatici e pulizia delle configurazioni.\tabularnewline
\textbf{Validazione} & Testare il sistema con utenti e operatori reali. & Feedback reale, dataset multimodale, gestione operatori e verifica e-mail.\tabularnewline
\textbf{Produzione} & Trasformare il prototipo in servizio reale. & Cloud, sicurezza, privacy, SLA e integrazione con fonti live.\tabularnewline
\bottomrule
\end{tabular}
}
\end{table}

Infine, il modello proposto non è limitato esclusivamente al contesto dei Giochi del Mediterraneo. La piattaforma potrebbe essere riutilizzata come soluzione \textit{SaaS} per altri eventi, poiché la struttura del sistema è modulare e indipendente dal dominio applicativo dettato dalla knowledge base. Possibili ambiti di applicazione sono comuni, università, aeroporti, ospedali o grandi eventi culturali. In questi contesti cambierebbero la knowledge base e alcune regole operative, ma la struttura principale del sistema resterebbe valida.